\documentclass[11pt]{article}

\usepackage[a4paper,margin=1in]{geometry}
\usepackage{amsmath,amssymb,mathtools}
\usepackage{booktabs}
\usepackage{enumitem}
\usepackage{xcolor}
\usepackage{hyperref}

\hypersetup{
  colorlinks=true,
  linkcolor=blue!55!black,
  urlcolor=blue!55!black,
  citecolor=blue!55!black
}

\setlist[itemize]{leftmargin=1.5em}

\newcommand{\R}{\mathbb{R}}
\newcommand{\norm}[1]{\left\lVert #1 \right\rVert_2}
\newcommand{\frob}[1]{\left\lVert #1 \right\rVert_F}

\title{Sliding-Window Complex Rules and Simple Functional Profiles}
\author{Companion note for \texttt{ar\_process\_sliding\_window\_embedding.ipynb}}
\date{}

\begin{document}
\maketitle

\section{Purpose}

The notebook constructs a sequence of simplicial complexes
\[
  L_1,L_2,\ldots,L_{k_{\max}}
\]
from a simulated autoregressive time series. For each embedding scale \(k\),
the notebook first builds a graph using one of the selected construction rules,
then takes the flag complex of that graph. The plotted profile is a simple
functional
\[
  k \longmapsto F(L_k), \qquad k=1,\ldots,k_{\max}.
\]

The following basic functionals are kept in the current notebook:
\[
  f_0,\quad f_1,\quad f_2,\quad \beta_0,\quad \beta_1,
  \quad \operatorname{PC}.
\]
These are the entries of the face vector, the first two Betti numbers, and the
point complexity. The notebook lets the user choose two such functionals
\(F_1,F_2\) and two scalar weights \(a,b\), and plots
\[
  aF_1(L_k)+bF_2(L_k).
\]
It also builds a scale-similarity matrix from the same linear combination:
\[
  S_{k\ell}=aF_1(L_k\oplus L_\ell)+bF_2(L_k\oplus L_\ell),
\]
where \(\oplus\) is either union or intersection, and plots the principal
eigenvector of this matrix against \(k\).

\section{Time-Series and Lifting Notation}

The simulated autoregressive process has order \(p\):
\[
  x_t = c + \phi_1 x_{t-1} + \cdots + \phi_p x_{t-p} + \varepsilon_t.
\]
In the interactive notebook, the user selects \(p\), not the coefficient vector.
The coefficients are drawn reproducibly from the seed in partial-autocorrelation
coordinates. Since every partial autocorrelation lies in \((-1,1)\), the
resulting AR polynomial is stationary.

For a univariate series \(x_0,\ldots,x_{n-1}\), the forward sliding-window
embedding of dimension \(k\) is
\[
  Y_t^{(k)} =
  (x_t,x_{t+1},\ldots,x_{t+k-1})^\top \in \R^k.
\]
The series is split chronologically into training and test segments. For each
\(k \leq k_{\max}\), the training windows of length \(k_{\max}\) are split as
\[
  W_t =
  (x_t,\ldots,x_{t+k_{\max}-1})^\top
  =
  \begin{pmatrix}
    Y_t^{(k)}\\
    Z_t^{(k)}
  \end{pmatrix},
\]
where
\[
  Z_t^{(k)} = (x_{t+k},\ldots,x_{t+k_{\max}-1})^\top.
\]

The completion matrix \(\widehat A_k\) is fitted by least squares:
\[
  \widehat A_k
  \in
  \arg\min_A
  \sum_t \left\| Z_t^{(k)} - A Y_t^{(k)} \right\|_2^2.
\]
For \(k=k_{\max}\), no tail is missing, and \(\widehat A_{k_{\max}}\) is empty.

On the test segment, each \(k\)-window is lifted to the common ambient space
\(\R^{k_{\max}}\):
\[
  p_t^{(k)}
  =
  \begin{pmatrix}
    Y_t^{(k)}\\
    \widehat A_k Y_t^{(k)}
  \end{pmatrix}.
\]
The corresponding prediction-error vector is
\[
  e_t^{(k)}
  =
  \widehat A_k Y_t^{(k)} - Z_t^{(k)}.
\]

\section{Common Graph and Flag-Complex Setup}

Fix \(k\). Let \(m\) be the number of complete test windows of length
\(k_{\max}\). Each time index \(t=0,\ldots,m-1\) is a vertex:
\[
  V_k = \{0,1,\ldots,m-1\}.
\]

Before applying a construction rule, the notebook restricts attention to
temporal candidate edges:
\[
  E_k^{\mathrm{temp}}
  =
  \bigl\{\{s,t\}: 0\leq s<t\leq m-1,\; 1\leq |s-t|\leq k-1\bigr\}.
\]
For \(k=1\), this set is empty. For larger \(k\), a vertex can connect only to
nearby window-start times.

Each construction rule produces a retained edge set
\[
  E_k \subseteq E_k^{\mathrm{temp}}.
\]
The graph is \(H_k=(V_k,E_k)\). The simplicial complex is the flag complex
\[
  L_k = \operatorname{Flag}(H_k).
\]
Thus a finite vertex set \(\sigma\subseteq V_k\) is a simplex of \(L_k\) exactly
when every pair of vertices in \(\sigma\) is an edge of \(H_k\). In the notebook,
GUDHI expands the graph through dimension two.


\section{Complex Construction Rules}

The notebook now uses manually chosen thresholds. Let \(\tau_e\geq0\) be the
error threshold, let \(\tau_d\geq0\) be the lifted-cloud distance threshold,
and let \(\tau_c\geq0\) be the covariance threshold.

A general way to describe the construction is to put nonnegative weights on
the temporal candidate graph. For a weight function
\[
  w_k:E_k^{\mathrm{temp}}\to\R_{\geq 0},
\]
retain the thresholded edge set
\[
  F_k^\tau
  =
  \{e\in E_k^{\mathrm{temp}}:w_k(e)\leq \tau\}.
\]
This gives a subgraph
\[
  H_k=(V_k,F_k^\tau)
\]
and therefore the data-dependent subcomplex
\[
  L_k=\operatorname{Flag}(H_k).
\]
If one wants the sequence \(L_1,\ldots,L_{k_{\max}}\) to be a filtration under
the common temporal edge sets, the edge weights should satisfy
\[
  w_{k+1}(e)\leq w_k(e),
  \qquad
  e\in E_k^{\mathrm{temp}},
  \quad
  k=1,\ldots,k_{\max}-1.
\]
The notebook reports the complexes scale by scale; this monotonicity condition
is the mathematical condition to check when the selected rule is intended to
define a true filtration.

\subsection{Both Errors Small}

Let
\[
  r_t^{(k)}=\norm{e_t^{(k)}}
\]
be the completion-error norm at time \(t\). A temporal candidate edge
\(\{s,t\}\in E_k^{\mathrm{temp}}\) is retained when both endpoint errors are
small:
\[
  \norm{e_s^{(k)}} \leq \tau_e,
  \qquad
  \norm{e_t^{(k)}} \leq \tau_e.
\]
This rule connects nearby windows only when both windows are individually well
predicted by the completion map.

\subsection{Similar Errors}

This rule ignores the absolute size of each individual error and only compares
error-vector similarity. A temporal candidate edge is retained when
\[
  \norm{e_s^{(k)}-e_t^{(k)}} \leq \tau_e.
\]
The endpoint errors may be large; what matters is that the two error vectors are
close to each other.

\subsection{Lifted-Cloud Distance Graph}

This rule uses the lifted points \(p_t^{(k)}\in\R^{k_{\max}}\). For all temporal
candidate edges, compute
\[
  d_{s,t}^{(k)}=\norm{p_s^{(k)}-p_t^{(k)}}.
\]
The edge \(\{s,t\}\) is retained when
\[
  d_{s,t}^{(k)} \leq \tau_d.
\]
This is a local distance graph on the lifted test cloud, restricted to temporal
candidate edges.

\subsection{Window Covariance Similarity}

Assume the AR process has finite innovation variance. The covariance rules are
therefore not available for Cauchy innovations or for centered Pareto
innovations with shape parameter at most \(2\).

Let
\[
  C_{s,t}^{(k)}
  =
  \operatorname{Cov}\!\left[Y_s^{(k)},Y_t^{(k)}\right].
\]
Using the stationary autocovariance function
\[
  \gamma(h)=\operatorname{Cov}(X_{u+h},X_u),
\]
the entries of this matrix are
\[
  C_{s,t}^{(k)}
  =
  \begin{pmatrix}
    \gamma(s-t) & \gamma(s-t-1) & \cdots & \gamma(s-t-k+1)\\
    \gamma(s-t+1) & \gamma(s-t) & \cdots & \gamma(s-t-k+2)\\
    \vdots & \vdots & \ddots & \vdots\\
    \gamma(s-t+k-1) & \gamma(s-t+k-2) & \cdots & \gamma(s-t)
  \end{pmatrix}.
\]
The notebook computes \(\gamma\) from the AR coefficients by solving the
stationary companion-matrix covariance equation, which is equivalent to the
Yule-Walker equations.

The scalar covariance summary is
\[
  \rho_k(s,t)=\frob{C_{s,t}^{(k)}}^2.
\]
Because the graph convention retains edges with small weights, the notebook
uses the reciprocal similarity weight
\[
  w_k^\rho(s,t)=\frac{1}{\rho_k(s,t)+\varepsilon},
  \qquad \varepsilon=10^{-12}.
\]
Thus an edge is retained when
\[
  w_k^\rho(s,t)\leq \tau_c.
\]
Equivalently, this keeps temporally admissible pairs whose window covariance is
large enough relative to the chosen threshold.

\subsection{Lifted Covariance Similarity}

The lifted point is
\[
  p_t^{(k)}
  =
  \begin{pmatrix}
    I_k\\
    \widehat A_k
  \end{pmatrix}
  Y_t^{(k)}.
\]
Therefore the cross-covariance matrix of two lifted points is
\[
  D_{s,t}^{(k)}
  =
  \operatorname{Cov}\!\left[p_t^{(k)},p_s^{(k)}\right]
  =
  \begin{pmatrix}
    I_k\\
    \widehat A_k
  \end{pmatrix}
  C_{s,t}^{(k)}
  \begin{pmatrix}
    I_k & \widehat A_k^\top
  \end{pmatrix}.
\]
The scalar lifted-covariance summary is
\[
  \eta_k(s,t)=\frob{D_{s,t}^{(k)}}^2.
\]
The notebook again converts this similarity score into a threshold weight:
\[
  w_k^\eta(s,t)=\frac{1}{\eta_k(s,t)+\varepsilon}.
\]
The edge \(\{s,t\}\) is retained when
\[
  w_k^\eta(s,t)\leq \tau_c.
\]
This rule compares windows after the least-squares completion map has lifted
them into the common ambient space \(\R^{k_{\max}}\).

\subsection{Full Temporal Graph}

This rule retains every temporal candidate edge:
\[
  E_k = E_k^{\mathrm{temp}}.
\]
It is useful as a baseline because it shows what the flag complex looks like
when the only restriction is temporal proximity.

\section{Simple Functional Definitions}

Let \(K=L_k\) be one of the complexes. Let \(K_i\) be the set of \(i\)-simplices.
In the notebook, dimensions \(0,1,2\) correspond to vertices, edges, and
triangles.

\subsection{Face-Vector Entries}

The face-vector entries are
\[
  f_i(K)=|K_i|.
\]
The notebook keeps
\[
  f_0(K)=\text{number of vertices},\qquad
  f_1(K)=\text{number of edges},\qquad
  f_2(K)=\text{number of triangles}.
\]

\subsection{Betti Numbers}

The Betti number \(\beta_i(K)\) is the rank of the \(i\)-th homology group over
\(\mathbb{F}_2\):
\[
  \beta_i(K)=\operatorname{rank} H_i(K;\mathbb{F}_2).
\]
The notebook keeps:
\begin{itemize}
  \item \(\beta_0(K)\), the number of connected components.
  \item \(\beta_1(K)\), the number of independent one-dimensional cycles.
\end{itemize}

\subsection{Point Complexity}

Let \(d=\dim K\), the largest simplex dimension present in \(K\). For a vertex
\(x\in K_0\), the point complexity is
\[
  \operatorname{pc}(x)
  =
  d-\max\{\dim(\sigma):\sigma\in K,\; x\in\sigma\}.
\]
Thus a vertex has low point complexity when it belongs to a high-dimensional
simplex, and high point complexity when it is isolated relative to the maximal
dimension of the complex.

The notebook plots the total point complexity
\[
  \operatorname{PC}(K)
  =
  \sum_{x\in K_0}\operatorname{pc}(x).
\]
Because the complexes are expanded through dimension two, the computed values
use \(d\leq 2\) in the current implementation.

\subsection{Linear Combination of Two Functionals}

In the interactive panel, the user chooses two basic functionals
\[
  F_1,F_2\in
  \{f_0,f_1,f_2,\beta_0,\beta_1,\operatorname{PC}\}
\]
and two scalar weights \(a,b\in\mathbb{R}\). The plotted profile is
\[
  k\longmapsto aF_1(L_k)+bF_2(L_k).
\]
For example, choosing \(F_1=\beta_0\), \(F_2=f_1\), \(a=2\), and \(b=-1\)
plots
\[
  2\beta_0(L_k)-f_1(L_k).
\]

\section{Scale-Similarity Matrix}

The notebook also compares different embedding scales pairwise. It uses the
same two functionals \(F_1,F_2\) and scalar weights \(a,b\) chosen for the
profile plot. Choose an operation
\[
  \oplus\in\{\cup,\cap\}.
\]
The scale-similarity matrix is
\[
  S_{k\ell}
  =
  aF_1(L_k\oplus L_\ell)+bF_2(L_k\oplus L_\ell),
  \qquad 1\leq k,\ell\leq k_{\max}.
\]
The union \(L_k\cup L_\ell\) contains all simplices appearing in either
complex. The intersection \(L_k\cap L_\ell\) contains exactly the simplices
appearing in both complexes. Since all \(L_k\) are stored as simplicial
complexes over the same test-window vertex set, both operations again produce
simplicial complexes.

The matrix \(S\) is symmetric. The notebook computes an eigenvector
\[
  v=(v_1,\ldots,v_{k_{\max}})^\top
\]
corresponding to the largest eigenvalue of \(S\), and plots
\[
  k\longmapsto v_k.
\]
The sign of an eigenvector is arbitrary, so the notebook flips the sign when
necessary to make the sum of entries nonnegative.

\section{Summary Table}

\begin{center}
\begin{tabular}{@{}ll@{}}
\toprule
Dropdown name & Mathematical meaning \\
\midrule
\texttt{f0 (vertices)} & \(f_0(K)=|K_0|\) \\
\texttt{f1 (edges)} & \(f_1(K)=|K_1|\) \\
\texttt{f2 (triangles)} & \(f_2(K)=|K_2|\) \\
\texttt{beta0} & \(\beta_0(K)\) \\
\texttt{beta1} & \(\beta_1(K)\) \\
\texttt{point complexity} & \(\operatorname{PC}(K)=\sum_{x\in K_0}\operatorname{pc}(x)\) \\
\bottomrule
\end{tabular}
\end{center}

\end{document}
